\documentclass[conference,letterpaper,10pt]{IEEEtran}
\IEEEoverridecommandlockouts

\usepackage[T1]{fontenc}
\usepackage{mathptmx}
\usepackage{amsmath,amssymb,amsfonts}
\usepackage{amsthm}
\usepackage{bm}
\usepackage{cite}
\usepackage{booktabs}
\usepackage{graphicx}
\usepackage{microtype}
\usepackage{balance}
\usepackage[hidelinks]{hyperref}

\pagestyle{empty}

\theoremstyle{plain}
\newtheorem{theorem}{Theorem}
\newtheorem{proposition}{Proposition}
\newtheorem{lemma}{Lemma}
\newtheorem{corollary}{Corollary}
\theoremstyle{remark}
\newtheorem{remark}{Remark}

\newcommand{\Nt}{N^{\ast}}
\newcommand{\xt}{x^{\ast}}
\newcommand{\edot}{\dot{e}}
\newcommand{\mudot}{\dot{\mu}}

\begin{document}

\title{Jevons' Paradox as a Safety-Critical Control Problem:\\
Consistent Rebound Dynamics and the Feasibility Limits of Fiscal Feedback}

\author{\IEEEauthorblockN{Peter Zsoldos}
\IEEEauthorblockA{Department of Computer Science and Electrical Engineering\\
Tampere University, Tampere, Finland\\
\texttt{zsoldosp44@gmail.com}}%
}

\maketitle
\thispagestyle{empty}

\begin{abstract}
Efficiency gains can raise total resource use---Jevons' paradox---which
invites a control formulation in which a dynamic resource tax rejects
efficiency shocks as disturbances. We show that the natural two-state
encoding of this idea, carrying resource draw $N$ and output $Y$ as
independent states alongside the definition $e=Y/N$ and an exogenous
efficiency path, is over-determined, is not invariant to the units of
the efficiency index, and structurally cannot represent efficiency
savings, so backfire is assumed rather than derived. Imposing the
identity yields a scalar plant that is linear in the geodesic
coordinate $x=\ln N$ and affine in the log gross tax $v=\ln(1+\tau)$
with constant actuation gain. In this plant backfire holds \emph{iff}
the demand elasticity satisfies $\eta>1$, and efficiency lowers
consumption on impact before raising it. Two conclusions reverse: the
disturbance that matters is the efficiency \emph{level}, not its
velocity, which is stabilising; and first-order Taylor prediction of
delayed data yields a neutral-type loop with exact margin
$\kappa\tau_d<0.9417$, \emph{below} the $\pi/2$ margin of the
unpredicted loop. The main result is a feasibility obstruction:
autonomous demand growth alone needs a $75\%$ peak tax and is
controllable, whereas adding the rebound channel raises the requirement
to $223\%$ and makes barrier feasibility fail for $\eta>0.948$---%
stricter than the backfire threshold. Feedback cuts cumulative
overshoot by $68\%$ but cannot render the safe set invariant at
credible tax ceilings.
\end{abstract}

\begin{IEEEkeywords}
Control barrier functions, predictor feedback, neutral delay systems,
actuator saturation, resource economics, rebound effect.
\end{IEEEkeywords}

\section{Introduction}
William Stanley Jevons observed in 1865 that efficiency gains in coal
use increased, rather than reduced, aggregate coal consumption
\cite{jevons1865}. When the economy-wide rebound exceeds $100\%$ the
condition is called backfire \cite{khazzoom1980,brookes1990,
saunders2008}. Its magnitude is contested: Gillingham \emph{et
al.}~\cite{gillingham2016} find little support for economy-wide
backfire, while sectoral ICT studies report rebound factors of
$1.2$--$1.6$ \cite{sorrell2007,coroama2019}. A useful model must
therefore treat the rebound magnitude as a parameter and report how
conclusions depend on it---the approach taken here.

A natural idea is to treat resource policy as feedback control: make
the tax a function of the measured state, so efficiency shocks are
rejected as disturbances rather than absorbed as consumption growth.
We pursued that programme and report that it survives only after
substantial repair. The two-state encoding of Section~\ref{sec:defects}
is our own first formulation, implemented and calibrated before we
found the defects reported below; we present it rather than quietly
discarding it because the defects are ones any first attempt is likely
to reproduce, and the corrected model reverses two conclusions we had
drawn from the flawed one.

\emph{Related work.} Rebound quantification dates to the
Khazzoom--Brookes postulate \cite{khazzoom1980,brookes1990}, surveyed
in \cite{sorrell2009,gillingham2016}; computable general equilibrium
approaches are descriptive, not suited to algorithmic policy design.
Ecological economists argue for absolute caps \cite{rockstrom2009},
and the choice between a price and a quantity instrument is
Weitzman's classical problem \cite{weitzman1974}: quantity dominates
when marginal damage is steep relative to marginal
abatement---precisely the case at a hard threshold, of which our
feasibility result is a dynamic restatement. On the control side,
forward invariance of a safe set is the province of control barrier
functions \cite{ames2017}, not control Lyapunov functions
\cite{kokotovic2001}---often conflated, and separated in
Section~\ref{sec:control}. For delay compensation the lineage is the
dead-time compensator \cite{smith1957}, finite spectrum assignment
\cite{manitius1979} and predictor feedback \cite{krstic2009}; stability
of the resulting neutral equations needs the difference-operator
machinery of \cite{hale1993} and the eigenvalue methods of
\cite{michiels2007}.

\emph{Contributions.} (i) Three independent defects in the natural
two-state formulation (Section~\ref{sec:defects}). (ii) A consistent
scalar plant (Section~\ref{sec:plant}) in which backfire is derived,
not assumed, and which reproduces the observed Jevons time signature.
(iii) A corrected delay-margin theorem (Section~\ref{sec:delay}): the
first-order Taylor predictor's exact margin is \emph{smaller} than the
unpredicted loop's. (iv) A feasibility obstruction
(Sections~\ref{sec:control}--\ref{sec:results}): because the required
tax is convex in the log-consumption gap, the rebound channel is the
increment that crosses the actuator bound despite not being the larger
contributor. (v) Negative findings reported rather than suppressed
(Section~\ref{sec:neg}).

\section{Model}

\subsection{Three Defects of the Two-State Formulation}
\label{sec:defects}
Our first formulation was the one that suggests itself immediately: state
$\bm{x}=[N,Y]^{\!\top}$, efficiency $e=Y/N$ declared as a definition,
an exogenously prescribed efficiency path $e(t)$, and dynamics
\begin{equation}
\label{eq:v1}
    \dot{\bm{x}} = f(\bm{x}) + g(\bm{x})\tau + \bm{D}\edot(t),
\end{equation}
with logistic $f$, multiplicative $g$ and constant coupling $\bm{D}$.
We implemented and calibrated \eqref{eq:v1} against IEA data-centre
electricity projections \cite{iea2024}. It has three independent
defects.

\emph{(D1) Over-determination.} If $e=Y/N$ is an identity and $N,Y$
each carry an independent ODE, $e(t)$ is determined by the state and
cannot also be prescribed; the prescribed $e(t)$ and the realised
$Y(t)/N(t)$ diverge by up to $6.4\times$ over $40$~yr.
\emph{(D2) Units dependence.} The rebound enters as $D_N\edot$, not
invariant under $e\mapsto ce$: rescaling the efficiency
index---physically the same technology in different units---moves peak
consumption from $3.96$ to $7.43$ normalised units. A prediction
depending on the units of an index is not a prediction.
\emph{(D3) No savings channel.} With $D_N>0$ and $\edot>0$ every
efficiency gain strictly increases $N$; no parameter gives
$\partial N/\partial e<0$, so backfire is assumed in \eqref{eq:v1}
rather than found.

These defects are structural, not calibration issues; every
quantitative claim derived from \eqref{eq:v1} must be regarded as an
artefact. The plant below is a replacement, not a refinement, and all
numbers in Section~\ref{sec:results} are recomputed from scratch.

\subsection{A Consistent Rebound Plant}
\label{sec:plant}
We keep the object that is genuinely exogenous---technology---and
impose the identity rather than duplicating it. Let
$\mu(t):=\ln e(t)$ be log efficiency. Because $e\mapsto ce$ shifts
$\mu$ by a constant, any model depending on $\mu$ only through
differences and derivatives is scale-free, curing~D2. Let $Y$ denote
demand for the energy \emph{service}, with effective service price
$p_{\text{eff}}(t)=(1+\tau(t))p_0/e(t)$, so that efficiency lowers and
the tax raises the cost of a unit of service. With constant-elasticity
long-run demand and partial adjustment at rate $\lambda>0$,
\begin{equation}
\label{eq:demand}
\begin{aligned}
    \tfrac{d}{dt}\ln Y &= \lambda\big(\ln Y^{\ast}(t) - \ln Y(t)\big), \\
    \ln Y^{\ast}(t) &= \ln A_0 + \gamma t + \eta\,\mu(t) - \eta\, v(t),
\end{aligned}
\end{equation}
where $\eta>0$ is the demand (rebound) elasticity, $\gamma$ the
autonomous income growth rate of service demand, and
$v:=\ln(1+\tau)\in\mathcal{V}=[0,v_{\max}]$, $v_{\max}=\ln(1+\tau_{\max})$,
is the \emph{log gross tax}. Imposing $N=Y/e$, i.e.\
$x:=\ln N=\ln Y-\mu$, gives the scalar plant
\begin{equation}
\label{eq:plant}
    \dot{x} = \underbrace{\lambda\big[\ln A_0 + \gamma t + (\eta-1)\mu(t) - x\big]
              - \mudot(t)}_{=:F(x,t)} \;+\; G\,v ,
\end{equation}
with \emph{constant} actuation gain $G=-\lambda\eta<0$. We normalise
$\mu(0)=0$, which is free by the scale-invariance above, and set
$\ln A_0=\ln N(0)$ so that the untaxed system starts on its frozen-time
equilibrium; all boundary pressure in what follows is therefore
generated by $\gamma$ and $\mu$ rather than by an initial transient.

Three structural gains follow. The plant is exactly affine in the
control; the actuation gain is constant, so the singularity
$g_N(\bm{x})=0$ that afflicts \eqref{eq:v1} cannot occur and no
policy-irrelevance caveat is needed; and \eqref{eq:plant} is
\emph{linear} in $x$, so the delay analysis of
Section~\ref{sec:delay} is exact rather than a linearisation. The
accessibility distribution of \eqref{eq:v1} degenerates on a
codimension-one locus; once the identity is imposed the system is
one-dimensional with $G\neq0$ everywhere, so that apparent geometric
richness was an artefact of the spurious degree of freedom introduced
by~D1.

\subsection{Backfire is Derived, Not Assumed}
\begin{proposition}[Backfire condition]
\label{prop:backfire}
Let $\mu$ and $v$ be constant. The frozen-time equilibrium of
\eqref{eq:plant} satisfies
$x_{\rm ss}=\ln A_0+\gamma t+(\eta-1)\mu-\eta v$, so
$\partial x_{\rm ss}/\partial\mu=\eta-1$. A permanent efficiency
improvement therefore raises long-run consumption iff $\eta>1$, leaves
it unchanged iff $\eta=1$, and lowers it iff $\eta<1$.
\end{proposition}
\begin{proof}
Set $\dot x=0$ and $\mudot=0$ in \eqref{eq:plant}, solve for $x$ and
differentiate in $\mu$. Along the ramp $\gamma t$ the solution tracks
$x_{\rm ss}$ with the constant offset $-\gamma/\lambda$, which is
independent of $\mu$ and therefore does not affect
$\partial x_{\rm ss}/\partial\mu$.
\end{proof}

\begin{proposition}[Jevons time signature]
\label{prop:signature}
The instantaneous response of consumption growth to an efficiency
impulse is $\partial\dot x/\partial\mudot=-1<0$ for every $\eta$, while
the long-run response is $\eta-1$. Efficiency therefore always reduces
consumption on impact, and the rebound arrives with the demand
adjustment time constant $1/\lambda$.
\end{proposition}
\begin{proof}
Read the coefficient of $\mudot$ in \eqref{eq:plant} and combine with
Proposition~\ref{prop:backfire}.
\end{proof}

Proposition~\ref{prop:backfire} recovers the textbook rebound threshold
as a property of the dynamics rather than an input, curing~D3; the
quasi-static target $x_{\rm ss}$ carries the efficiency contribution
$(\eta-1)\mu(t)$ exactly, by construction, and we verified this
identity numerically across $\eta\in[0.4,1.8]$, where
$M:=\mu(T)-\mu(0)$. Proposition~\ref{prop:signature} is the content of
Fig.~\ref{fig:mechanism}(a), which plots the \emph{realised} trajectory
rather than this quasi-static target: for $\eta=1.4$ the efficiency
contribution dips to $-0.42$ near $t=13$~yr, then recovers to only
$+0.68$ by the horizon $t=25$~yr, short of the quasi-static
$(\eta-1)M=0.88$ by the demand-adjustment lag $1/\lambda=4$~yr---so the
$25$-year window in Table~\ref{tab:calib} understates the eventual
backfire rather than overstating it. The two-state model \eqref{eq:v1}
cannot produce this sign reversal at any parameter value.
Fig.~\ref{fig:mechanism}(b) sweeps
$\eta$ and locates the two thresholds that organise the paper: the
backfire condition $\eta=1$ and the feasibility limit
$\eta_{\rm crit}=0.948$ of Corollary~\ref{cor:etacrit}.

\begin{figure}[!t]
    \centering
    \includegraphics[width=\columnwidth]{fig_mechanism}
    \caption{Corrected rebound mechanism. (a) The efficiency
    contribution to $\ln N$ is negative on impact and positive in the
    long run, reversing sign near $t=17.6$~yr for $\eta=1.4$; the
    two-state model cannot produce this. (b) Long-run consumption
    versus rebound elasticity, with the derived backfire threshold
    $\eta=1$ and the feasibility limit $\eta_{\rm crit}=0.948$.}
    \label{fig:mechanism}
\end{figure}

\begin{remark}[The disturbance is the level, not the velocity]
\label{rem:level}
In \eqref{eq:plant} the efficiency level enters with coefficient
$\lambda(\eta-1)$ and its velocity with coefficient $-1$. The velocity
term is therefore \emph{stabilising}: rapid innovation transiently
relieves rather than aggravates boundary pressure. This reverses the
premise that $\edot$ is the disturbance to be rejected. The quantity
that must be tracked is the accumulated level $(\eta-1)\mu(t)$.
\end{remark}

\begin{table}[!t]
\centering
\caption{Calibration ($N$ normalised by $\Nt=460$~TWh/yr)}
\label{tab:calib}
\setlength{\tabcolsep}{3.5pt}
\begin{tabular}{llcl}
\toprule
Symbol & Meaning & Value & Basis \\
\midrule
$\Nt$        & Resource ceiling            & $1.0$    & IEA high scenario \cite{iea2024} \\
$N(0)$       & 2023 data-centre draw       & $0.804$  & IEA \cite{iea2024} \\
$A_0$        & Demand scale                & $0.804$  & $x(0)$ on frozen-time eq. \\
$\mu(0)$     & Efficiency index origin     & $0$      & normalisation (D2) \\
$\lambda$    & Demand adjustment rate      & $0.25$   & $\approx2.8$~yr half-life \\
$\gamma$     & Autonomous service growth   & $0.04$   & long-run digital demand \\
$\eta$       & Rebound elasticity          & $1.40$   & ICT range \cite{sorrell2007,coroama2019} \\
$M$          & Total log efficiency gain   & $\ln 9$  & three GPU generations \\
$\tau_{\max}$& Actuator ceiling            & $1.0$    & $100\%$ tax \\
$\tau_d$     & Reporting delay             & $1.0$~yr & macro release lag \\
$K$          & Lyapunov gain               & $2.0$    & design choice \\
$T$          & Horizon                     & $25$~yr  & 2025--2050 \\
\bottomrule
\end{tabular}
\end{table}

\section{Geometry, Control Design and Feasibility}
\label{sec:control}
Let the safe set be $\Omega=\{x:h(x)\ge0\}$ with $h(x):=\xt-x$ and
$\xt=\ln\Nt$, and take the quadratic Lyapunov function
$V(x)=\tfrac12(x-\xt)^2$ in the log coordinate: $x=\ln N$ is where the
plant \eqref{eq:plant} is linear, and it is also the geodesic
coordinate of the metric $g_{NN}=1/(\sigma N^2)$ for an elasticity of
substitution $\sigma$, which enters $V$ only as a positive constant and
so cancels from $\dot V\le-KV$---useful, since $\sigma$ is among the
least reliably identified parameters in applied production theory.
Also $V\to\infty$ as $N\to0^{+}$, so positivity of the resource state
is preserved structurally.

\subsection{Lyapunov Tracking and Barrier Filtering}
Imposing $\dot V\le-KV$ on \eqref{eq:plant} and writing $z:=x-\xt$
gives, in the interior of $\mathcal{V}$,
\begin{equation}
\label{eq:clf}
    v_{\rm CLF} = \left[\frac{F(x,t) + \kappa\,z}{-G}\right]_0^{v_{\max}},
    \qquad \kappa := \frac{K}{2},
\end{equation}
where $[\cdot]_0^{v_{\max}}$ denotes clipping to $\mathcal{V}$; the
unsaturated closed loop is $\dot z=-\kappa z$. The factor $\kappa=K/2$
follows from $\dot V\le-KV$ with $V=\tfrac12z^2$; all delay conditions
below are stated in $\kappa$, and expressing them in $K$ introduces a
factor-of-two error.

A Lyapunov function certifies convergence to a target; it does not
certify that a boundary is never crossed. With class-$\mathcal{K}$ gain
$\alpha>0$, the invariance condition $\dot h\ge-\alpha h$ gives, since
$G<0$,
\begin{equation}
\label{eq:cbf}
    v \;\ge\; v_{\rm CBF}(x,t) := \frac{F(x,t)-\alpha\,(\xt-x)}{-G},
\end{equation}
and we deploy the minimum-intervention filter
$v^{\star}=[v_{\rm CBF}(x,t)]_0^{v_{\max}}$. Far below the ceiling $h$
is large and \eqref{eq:cbf} is slack, so $v_{\rm CBF}<0$ and no tax is
levied: the nominal policy is no tax at all and the barrier is the only
thing that ever raises $\tau$. If $x(0)\in\Omega$ and
$v_{\rm CBF}(x(t),t)\le v_{\max}$ for all $t\ge0$, then $v^\star$
satisfies $\dot h\ge-\alpha h$ pointwise, so
$h(t)\ge h(0)e^{-\alpha t}\ge0$ by the comparison lemma and $\Omega$ is
forward invariant \cite{ames2017}.

\subsection{The Feasibility Obstruction}
\label{sec:feas}
Evaluating \eqref{eq:cbf} on $\partial\Omega$ defines the tax required
to hold the ceiling,
\begin{equation}
\label{eq:vreq}
    v_{\rm req}(t) = \frac{F(\xt,t)}{-G}
    = \frac{\lambda\big[\ln A_0+\gamma t+(\eta-1)\mu(t)-\xt\big]-\mudot(t)}
           {\lambda\eta},
\end{equation}
with $\tau_{\rm req}=e^{v_{\rm req}}-1$.

\begin{proposition}[Admissibility]
\label{prop:feas}
$\Omega$ admits an invariance certificate at time $t$ iff
$v_{\rm req}(t)\le v_{\max}$. Under Table~\ref{tab:calib} this fails on
$30.8\%$ of the horizon, first at $t=17.3$~yr, with
$\max_t\tau_{\rm req}=223\%$. Setting $M=0$ (income growth only) gives
$\max_t\tau_{\rm req}=75\%$ and the condition never fails.
\end{proposition}

The decomposition behind Proposition~\ref{prop:feas} is the paper's
central quantitative point. At the horizon the log-consumption gap the
tax must close is
$\underbrace{\gamma T}_{1.000\,(53\%)}+\underbrace{(\eta-1)M}_{0.879\,(47\%)}+\ln A_0$,
so the two drivers are of comparable size. But $v_{\rm req}$ is linear
in the gap while $\tau_{\rm req}=e^{v_{\rm req}}-1$ is exponential in
it: a gap of $0.78$ needs $75\%$, a gap of $1.66$ needs $223\%$. The
rebound channel is not the larger contributor, yet it is the increment
that carries the requirement across the actuator bound, because the
marginal cost of closing the gap rises exponentially
(Fig.~\ref{fig:feasibility}).

\begin{figure}[!t]
    \centering
    \includegraphics[width=\columnwidth]{fig_feasibility}
    \caption{The central result. Income growth alone requires a peak
    tax of $75\%$ and never exceeds the credible ceiling. Adding the
    rebound channel raises the requirement to $223\%$ and makes the
    invariance certificate infeasible from $t=17.3$~yr. The early dip
    below zero is the transient savings channel of
    Proposition~\ref{prop:signature}.}
    \label{fig:feasibility}
\end{figure}

\begin{remark}[Convexity of the actuation requirement]
\label{rem:convex}
The mechanism above is convexity, not arithmetic, with a consequence
for how the requirement should be reported. Write
$\Gamma(t):=\ln A_0+\gamma t+(\eta-1)\mu(t)-\xt$ for the
log-consumption gap, so $v_{\rm req}$ is affine in $\Gamma$ while
$\tau_{\rm req}=e^{v_{\rm req}}-1$ is strictly convex in it. So
attributing infeasibility to a single channel is meaningful only at
the margin, as in Proposition~\ref{prop:feas}; and if $\Gamma$ is
uncertain, Jensen's inequality gives
$\mathbb{E}[\tau_{\rm req}]\ge e^{\mathbb{E}[v_{\rm req}]}-1$, with
equality only for a degenerate distribution---the tax required on
average strictly exceeds the tax at the nominal calibration, and the
gap widens with the dispersion of $(\eta,\gamma,M)$. An actuator
ceiling sized at nominal values is therefore systematically
under-sized; this argues for reporting the distribution of
$\tau_{\rm req}$ rather than a point value.
\end{remark}

\begin{corollary}[Critical elasticity]
\label{cor:etacrit}
At $\tau_{\max}=1$ the condition of Proposition~\ref{prop:feas} holds
for all $t$ iff $\eta\le\eta_{\rm crit}=0.948$. Since
$\eta_{\rm crit}<1$, controllability is lost \emph{before} backfire
sets in: the policy-relevant threshold is stricter than the economic
one.
\end{corollary}
\begin{proof}
Since $\eta$ enters \eqref{eq:vreq} both as a multiplier of $\mu(t)$
and as the divisor, monotonicity in $\eta$ is not immediate and
deserves an argument rather than an appeal to Fig.~\ref{fig:mechanism}(b)
alone. Write
\begin{equation}
\label{eq:vreq-split}
\begin{aligned}
    v_{\rm req}(\eta,t) &= \mu(t) + \frac{\Psi(t)}{\lambda\eta}, \\
    \Psi(t) &:= \lambda\big[\ln A_0+\gamma t-\xt-\mu(t)\big] - \mudot(t),
\end{aligned}
\end{equation}
obtained by expanding \eqref{eq:vreq}; $\Psi$ does not depend on
$\eta$. Hence for fixed $t$, $v_{\rm req}(\cdot,t)$ is monotonic on
$(0,\infty)$ with $\partial v_{\rm req}/\partial\eta=-\Psi(t)/(\lambda\eta^2)$,
strict whenever $\Psi(t)\neq0$. As the efficiency path saturates,
$\mudot(t),\ddot\mu(t)\to0$ near $t=T$, so
$\partial v_{\rm req}/\partial t=\mudot(t)+\dot\Psi(t)/(\lambda\eta)
\to\gamma/\eta>0$ there for every $\eta>0$: the requirement is
eventually increasing in $t$ regardless of elasticity. Under
Table~\ref{tab:calib} this terminal regime governs the maximum,
$g(\eta):=\max_{t\in[0,T]}v_{\rm req}(\eta,t)=v_{\rm req}(\eta,T)
=\mu(T)+\Psi(T)/(\lambda\eta)$, throughout $\eta\in[0.7,2.0]\ni\eta_{\rm crit}$,
which we confirm numerically; for smaller $\eta$ the interior maximum
occurs during the transient dip of Proposition~\ref{prop:signature}
and is negative, hence trivially below $v_{\max}$. Since
$\Psi(T)<0$ at this calibration, $g$ is strictly increasing on
$[0.7,2.0]$, so it crosses $v_{\max}$ exactly once; $\eta_{\rm crit}$
is that crossing, and $g(\eta)\le v_{\max}$ holds iff
$\eta\le\eta_{\rm crit}$.
\end{proof}

\begin{lemma}[ISS with respect to estimate error]
\label{lem:iss}
Let the terms of $F$ be estimated with total error bounded by $\delta$
and let the actuator be unsaturated. Then
$|z(t)|\le|z(0)|e^{-\kappa t}+\delta/\kappa$.
\end{lemma}
\begin{proof}
The perturbed loop is $\dot z=-\kappa z+w$ with $|w|\le\delta$; apply
the comparison lemma to $\tfrac{d}{dt}|z|\le-\kappa|z|+\delta$.
\end{proof}

Lemma~\ref{lem:iss} holds only while $v\in(0,v_{\max})$. Once
$v\equiv v_{\max}$ the loop is open and no input-to-state bound
applies---a caveat binding for roughly $30\%$ of our horizon.

\section{Delay Compensation}
\label{sec:delay}
Macroeconomic aggregates are released with a lag
$\tau_d\approx0.5$--$1.5$~yr. The natural first attempt extrapolates
one delay interval forward,
$\hat{x}(t)=x(t-\tau_d)+\tau_d\dot{x}(t-\tau_d)$, giving the
closed-loop error equation
\begin{equation}
\label{eq:neutral}
    \dot{z}(t) = -\kappa\big[z(t-\tau_d) + \tau_d\,\dot{z}(t-\tau_d)\big],
\end{equation}
which contains the delayed \emph{derivative} and is therefore of
neutral type \cite{hale1993}. Equivalently
$\frac{d}{dt}[z(t)+\kappa\tau_d z(t-\tau_d)]=-\kappa z(t-\tau_d)$,
whose difference operator is stable iff $\kappa\tau_d<1$. Since
\eqref{eq:plant} is exactly linear, \eqref{eq:neutral} is the exact
closed loop.

\begin{theorem}[Exact delay margin of the Taylor predictor]
\label{thm:taylor}
Let $\kappa>0$, $\tau_d>0$ and $h:=\kappa\tau_d$. The loop
\eqref{eq:neutral} is exponentially stable iff
\begin{equation}
    h < h^{\ast} = \frac{\nu^{\ast}}{\sin\nu^{\ast}\big(1+(\nu^{\ast})^{2}\big)}
    \approx 0.941680,
\end{equation}
where $\nu^{\ast}\approx2.798386$ is the smallest positive root of
$\cos\nu+\nu\sin\nu=0$.
\end{theorem}
\begin{proof}
The characteristic function is
$\Delta(s)=s+\kappa(1+\tau_d s)e^{-s\tau_d}$. With $\lambda_c=s\tau_d$
and $h=\kappa\tau_d$ this is
$\lambda_c+h(1+\lambda_c)e^{-\lambda_c}=0$. As $h\to0^{+}$ the finite
root is $\lambda_c\approx-h<0$, so the loop begins stable and
instability requires an imaginary-axis crossing. Put
$\lambda_c=j\nu$, $\nu>0$. Using
$(1+j\nu)(\cos\nu-j\sin\nu)=(\cos\nu+\nu\sin\nu)+j(\nu\cos\nu-\sin\nu)$,
\begin{align}
    \text{Re}:&\quad h(\cos\nu+\nu\sin\nu)=0, \label{eq:re}\\
    \text{Im}:&\quad \nu+h(\nu\cos\nu-\sin\nu)=0. \label{eq:im}
\end{align}
As $h>0$, \eqref{eq:re} forces $\cos\nu=-\nu\sin\nu$; substituting into
\eqref{eq:im} gives $h=\nu/(\sin\nu-\nu\cos\nu)=\nu/[\sin\nu(1+\nu^{2})]$.
The smallest positive root of $\cos\nu+\nu\sin\nu=0$ is $\nu^\ast$,
giving $h^\ast$; larger roots give larger $h$ and so do not govern the
first crossing. Since $h^\ast<1$ the difference operator is stable
throughout $[0,h^\ast)$, no root chain escapes to $+\infty$, and the
root condition is necessary and sufficient \cite{hale1993,michiels2007}.
\end{proof}

\begin{proposition}[No prediction beats first-order prediction]
\label{prop:nopred}
Feeding back the delayed measurement without prediction gives
$\dot z(t)=-\kappa z(t-\tau_d)$, exponentially stable iff
$\kappa\tau_d<\pi/2$. The Taylor predictor therefore \emph{reduces} the
delay margin by the factor $\pi/(2h^{\ast})\approx1.668$, a loss of
about $40\%$.
\end{proposition}
\begin{proof}
For $\Delta(s)=s+\kappa e^{-s\tau_d}$ and $s=j\omega$ the real part
gives $\cos\omega\tau_d=0$, so $\omega\tau_d=\pi/2$, and the imaginary
part gives $\omega=\kappa$. Compare with Theorem~\ref{thm:taylor}.
\end{proof}

Both constants were confirmed by numerical continuation of the
rightmost characteristic root; the computed sign changes bracket
$0.9417$ and $1.5708$ (Fig.~\ref{fig:delay}(a));
Fig.~\ref{fig:delay}(b) shows the corresponding unsaturated
trajectories on either side of $h^\ast$.

\begin{figure}[!t]
    \centering
    \includegraphics[width=\columnwidth]{fig_delay}
    \caption{Corrected delay analysis. (a) Rightmost characteristic
    root versus $h=\kappa\tau_d$: the first-order Taylor predictor
    loses stability at $h^\ast\approx0.942$, \emph{before} the
    unpredicted loop at $\pi/2$. (b) Simulation of the exactly linear
    plant without saturation.}
    \label{fig:delay}
\end{figure}

\begin{remark}[Why a retarded Lyapunov--Krasovskii argument fails]
\label{rem:badproof}
It is tempting---and we initially attempted---to certify
\eqref{eq:neutral} with
$V_{LK}=\tfrac12z^2(t)+\tfrac{\kappa\tau_d}{2}\int_{t-\tau_d}^{t}z^2(\theta)\,d\theta$
and conclude $\tau_d<1/\kappa$. This fails, and worth recording because
the bound errs in \emph{direction}, not merely tightness: a memory
term penalising $z^2$ cannot dominate the cross term in
$\dot z(t-\tau_d)$ that a \emph{neutral} equation produces---a
functional penalising $\dot z^2$ over the window is required
\cite{hale1993}---and closing the argument by asserting
$z^2(t-\tau_d)\ge\tau_d^2\dot z^2(t-\tau_d)$ is circular, since along
$\dot z=-\kappa z$ it reads $1-(\kappa\tau_d)^2\ge0$, the conclusion
itself. The resulting $\kappa\tau_d<1$ coincides with the (necessary
but weaker) difference-operator condition, not $h^\ast$. A linear
delay equation does not undergo a Hopf bifurcation; it exhibits a root
crossing.
\end{remark}

The remedy is the classical dead-time compensator
\cite{smith1957,manitius1979,krstic2009}: integrate the model across
the unmeasured window rather than extrapolating linearly,
\begin{equation}
\label{eq:smith}
    \hat{x}(t) = x(t-\tau_d) + \int_{t-\tau_d}^{t}\!\big[F(\hat{x},s) + G\,v(s)\big]\,ds ,
\end{equation}
evaluated by fixed-step RK4 using the known applied input history.
Under an exact model the prediction is exact, the delay is cancelled
and $\dot z=-\kappa z$ is recovered for arbitrary $\tau_d$. An
innovation term $L(x(s)-\hat{x}(s))$ may be appended to
\eqref{eq:smith} only for $s\le t-\tau_d$; inside the prediction window
no measurement exists, so the integration there must be open loop.
Implementations that interpolate the true state across the window are
non-causal. Enforcing strict causality changes reported metrics by less
than $0.2\%$ under a well-matched model, but the distinction is
required for correctness and would matter under larger mismatch.

\section{Numerical Results}
\label{sec:results}
Equation~\eqref{eq:plant} was integrated with adaptive RK45
($\mathrm{rtol}=10^{-7}$, $\mathrm{atol}=10^{-9}$) over $T=25$~yr under
Table~\ref{tab:calib}. Metrics are peak $N$; integrated absolute error
$\mathrm{IAE}=\int|N-\Nt|\,dt$; integrated squared effort
$\mathrm{ISCE}=\int\tau^2dt$; terminal output $Y_f$; saturation duty
$S$; and whether the trajectory breaches $1.05\Nt$.

\begin{table}[!t]
\centering
\caption{Controller performance ($\eta=1.4$, $\tau_d=1$~yr, $\tau_{\max}=1$, $T=25$~yr)}
\label{tab:metrics}
\setlength{\tabcolsep}{3.5pt}
\begin{tabular}{lcccccc}
\toprule
Controller & peak $N$ & IAE & ISCE & $Y_f$ & $S$ [\%] & breach \\
\midrule
Open loop                  & $3.682$ & $14.30$ & $0.00$ & $32.93$ & $0.0$  & yes \\
Static tax $22\%$          & $2.789$ & $12.09$ & $1.21$ & $24.94$ & $0.0$  & yes \\
CLF, no lag                & $1.543$ & $4.64$  & $8.55$ & $13.80$ & $30.3$ & yes \\
CLF $+$ Taylor pred.       & $1.542$ & $4.65$  & $8.55$ & $13.79$ & $29.6$ & yes \\
CLF $+$ Smith pred.        & $1.542$ & $\bm{4.63}$ & $8.56$ & $13.79$ & $30.3$ & yes \\
CBF filter $+$ Smith pred. & $\bm{1.538}$ & $4.65$ & $8.59$ & $13.76$ & $30.0$ & yes \\
\bottomrule
\end{tabular}
\end{table}

\emph{Attenuation without invariance.} Table~\ref{tab:metrics} gives
the primary comparison: closed-loop feedback reduces IAE from $14.30$
to $4.63$ ($67.6\%$) and peak consumption from $3.682\Nt$ to
$1.542\Nt$ ($58.1\%$)---both large and policy-relevant, but neither is
invariance: every admissible controller breaches the ceiling,
consistent with Proposition~\ref{prop:feas}, with the actuator
saturated for about $30\%$ of the horizon (Fig.~\ref{fig:controllers}).

\emph{Predictor comparison.} At the nominal calibration
$\kappa\tau_d=1.0$, only marginally beyond $h^\ast$, the two predictors
are close in Table~\ref{tab:metrics}; Table~\ref{tab:delay} separates
them in the unsaturated regime: the Taylor predictor's residual
oscillation grows by a factor of $768$ across the sweep, inflecting
sharply at $\kappa\tau_d\approx h^\ast$ as Theorem~\ref{thm:taylor}
predicts, against $4.6$ for the model-based predictor---inverting the
intuition, which we initially followed, that first-order prediction
enlarges the stability region.

\begin{table}[!t]
\centering
\caption{Residual oscillation of $N$ vs.\ reporting delay ($\kappa=1$, unsaturated)}
\label{tab:delay}
\setlength{\tabcolsep}{6pt}
\begin{tabular}{cccc}
\toprule
$\tau_d$ [yr] & $\kappa\tau_d$ & Taylor pred. & Smith pred. \\
\midrule
$0.50$ & $0.50$ & $0.0011$ & $0.0008$ \\
$0.90$ & $0.90$ & $0.0034$ & $0.0013$ \\
$1.00$ & $1.00$ & $0.0066$ & $0.0014$ \\
$1.25$ & $1.25$ & $0.0439$ & $0.0019$ \\
$1.50$ & $1.50$ & $0.3451$ & $0.0024$ \\
$2.00$ & $2.00$ & $0.8310$ & $0.0038$ \\
\bottomrule
\end{tabular}
\end{table}

\begin{figure}[!t]
    \centering
    \includegraphics[width=\columnwidth]{fig_controllers}
    \caption{Trajectories at $\tau_d=1$~yr. (a) Resource draw: feedback
    delays and flattens the breach but does not prevent it. (b) Fiscal
    signal: the tax is railed at $\tau_{\max}$ over the final third of
    the horizon, which is why controller-design choices have little
    leverage.}
    \label{fig:controllers}
\end{figure}

\emph{Required authority and its cost.} Sweeping the actuator ceiling,
invariance first obtains at $\tau_{\max}\approx220\%$---close to the
analytic $223\%$ of Proposition~\ref{prop:feas}---at a cost of $35.0\%$
of terminal output ($13.76\to8.94$): defending the $460$~TWh/yr
ceiling against these paths costs roughly a third of terminal digital
output.

\emph{Robustness.} Over $300$ Monte Carlo draws of the four genuinely
uncertain quantities $(\eta,\lambda,\gamma,M)$, the open loop breaches
in $94\%$ of draws and the closed loop in $77\%$; no draw destabilises,
but the residual breach rate is governed by whether $\eta>1$, not
control design. Consistent with Remark~\ref{rem:convex}, the same
draws give a mean peak requirement
$\mathbb{E}[\max_t\tau_{\rm req}]=241\%$ against the nominal $223\%$
and a 90th percentile of $318\%$: sizing $\tau_{\max}$ from the nominal
calibration understates the requirement at every quantile above the
median.

\subsection{Negative Findings}
\label{sec:neg}
Several comparisons one might expect to favour the proposed machinery
do not: with a non-negative tax the barrier and Lyapunov designs
coincide almost exactly (peak $1.538$ vs.\ $1.542$; $Y_f$ $13.76$
vs.\ $13.79$) and, with ample authority, both converge onto the ceiling
($N=0.9994$ and $0.9996$)---we had expected the barrier to preserve a
margin, but the disturbance drives $N$ upward and a one-sided tax can
only resist. Varying the class-$\mathcal{K}$ gain $\alpha$ over
$[0.2,2.0]$ moves peak $N$ only from $1.519$ to $1.544$, and holding
the controller's internal model at nominal while perturbing the plant
changes the $95$th-percentile peak from $1.887$ to $1.890$. All three
share one explanation: for roughly a third of the horizon the actuator
is railed, and a railed actuator consults neither its model, its gain,
nor its certificate---itself a reason to report the feasibility
boundary rather than a controller ranking. The barrier's value here is
the \emph{certificate} \eqref{eq:vreq}, not improved trajectories.

\section{Discussion and Conclusion}
Treating resource policy as feedback is well posed and worthwhile: it
delivers a $68\%$ reduction in cumulative boundary violation and a
$58\%$ reduction in peak draw from a signal computable from delayed
data, provided the predictor is model-based rather than a truncated
expansion. What it does not buy is the elimination of backfire. The
obstruction is structural: because the required tax is exponential in
the log-consumption gap, a price instrument capable of defending a
hard ceiling against sustained demand growth plus rebound must be so
large that its political and distributional feasibility becomes the
binding constraint instead---the dynamic counterpart of Weitzman's
result \cite{weitzman1974} that, near a steep threshold, quantity
instruments dominate price instruments. The constructive reading is
that $v_{\rm req}-v_{\max}$ in \eqref{eq:vreq} is computable in real
time from observable data and is precisely the signal indicating when
a price instrument must be supplemented by a quantity instrument.

Two reversals relative to the velocity-rejection view deserve emphasis:
by Remark~\ref{rem:level} the disturbance to be tracked is the
efficiency \emph{level}, not its velocity, so a controller built to
reject $\edot$ is rejecting a stabilising signal; and by
Proposition~\ref{prop:nopred} the intuitive delay fix makes matters
worse. Both errors are invisible in a model with defects D1--D3---the
practical argument for dimensional and definitional consistency before
control design.

\emph{Limitations.} Equation~\eqref{eq:plant} is a scalar reduced form
with one instrument---no capital stock, sectoral disaggregation, trade
leakage, endogenous innovation response, or forward-looking
expectations. The last is most serious: an agent anticipating a rising
tax would reallocate intertemporally, a channel a partial-adjustment
rule cannot represent. The headline depends on $\eta>1$, which the
economy-wide literature does not establish \cite{gillingham2016};
Corollary~\ref{cor:etacrit} and the full sweep, not a point claim, are
the appropriate summary. We model a pure reporting delay rather than
quarterly sampling, which would add roughly half a sampling interval of
effective delay and require discrete-time barrier conditions; $Y_f$ is
a scalar output proxy with no welfare or incidence analysis. Future
work should replace the pointwise barrier with a receding-horizon
formulation \cite{mayne2014} over a forecast $\mu$ trajectory, treat
$\tau_{\max}$ as a decision variable in a joint price--quantity design,
and identify $\eta$ and $\lambda$ from data with quantified
uncertainty, since the policy conclusion turns on which side of
$0.948$ the elasticity lies.

Jevons' paradox is a legitimate object of control-theoretic study, but
the obvious state-space encoding of it is not. Imposing the identity
yields a scalar plant in which backfire holds iff $\eta>1$ and
efficiency gains lower consumption before raising them. In that plant
delay compensation by first-order Taylor extrapolation \emph{shrinks}
the delay margin from $\kappa\tau_d<\pi/2$ to $\kappa\tau_d<0.9417$.
Autonomous demand growth alone is controllable with a $75\%$ tax;
adding the rebound channel raises the requirement to $223\%$, and the
feasibility limit $\eta_{\rm crit}=0.948$ is stricter than the backfire
threshold. Fiscal feedback is therefore necessary for boundary defence,
demonstrably insufficient on its own, and the safety-critical
formulation makes that insufficiency measurable in advance.

\section*{Acknowledgment}
Portions of this manuscript were drafted with the assistance of a large
language model (Anthropic Claude), which was used for prose drafting
and \LaTeX{} preparation in
Sections~\ref{sec:defects}--\ref{sec:neg}; the model formulation,
proofs, numerical results and all conclusions were derived and verified
by the author, who takes full responsibility for the content.

\section*{Reproducibility}
Code and calibration for every figure and table are at
\url{https://github.com/Hannibal0319/jevons-paradox-control}:
\texttt{tables.py} reproduces Tables~\ref{tab:metrics}--\ref{tab:delay}
to the digit, \texttt{figures.py} reproduces
Figs.~\ref{fig:mechanism}--\ref{fig:controllers}.
Propositions~\ref{prop:backfire}--\ref{prop:signature} and
Corollary~\ref{cor:etacrit} were checked against simulation; the delay
margins of Theorem~\ref{thm:taylor} and Proposition~\ref{prop:nopred}
against numerical continuation of the rightmost characteristic root.

\balance

\begin{thebibliography}{99}

\bibitem{jevons1865}
W.~S. Jevons, \emph{The Coal Question}.\hskip 1em plus 0.5em minus 0.4em\relax
London, UK: Macmillan, 1865.

\bibitem{khazzoom1980}
J.~D. Khazzoom, ``Economic implications of mandated efficiency in standards for
household appliances,'' \emph{Energy J.}, vol.~1, no.~4, pp. 21--40, 1980.

\bibitem{brookes1990}
L.~Brookes, ``Energy efficiency and economic growth,'' \emph{Energy Policy},
vol.~18, no.~2, pp. 199--201, 1990.

\bibitem{saunders2008}
H.~D. Saunders, ``Fuel conserving (and using) production functions,''
\emph{Energy Econ.}, vol.~30, no.~5, pp. 2184--2235, 2008.

\bibitem{sorrell2009}
S.~Sorrell, ``Jevons' paradox revisited: The evidence for backfire from improved
energy efficiency,'' \emph{Energy Policy}, vol.~37, no.~4, pp. 1456--1469, 2009.

\bibitem{sorrell2007}
S.~Sorrell, \emph{The Rebound Effect}.\hskip 1em plus 0.5em minus 0.4em\relax
London, UK: UK Energy Research Centre, 2007.

\bibitem{gillingham2016}
K.~Gillingham, D.~Rapson, and G.~Wagner, ``The rebound effect and energy
efficiency policy,'' \emph{Rev. Environ. Econ. Policy}, vol.~10, no.~1,
pp. 68--88, 2016.

\bibitem{coroama2019}
V.~C. Coroam\u{a} and F.~Mattern, ``Digital rebound: Why digitalization will not
redeem us our environmental sins,'' in \emph{Proc. ICT4S}, Lappeenranta,
Finland, 2019.

\bibitem{rockstrom2009}
J.~Rockstr{\"o}m \emph{et al.}, ``A safe operating space for humanity,''
\emph{Nature}, vol.~461, no.~7263, pp. 472--475, 2009.

\bibitem{weitzman1974}
M.~L. Weitzman, ``Prices vs. quantities,'' \emph{Rev. Econ. Stud.}, vol.~41,
no.~4, pp. 477--491, 1974.

\bibitem{kokotovic2001}
P.~Kokotovi\'{c} and M.~Arcak, ``Constructive nonlinear control: A historical
perspective,'' \emph{Automatica}, vol.~37, no.~5, pp. 637--662, 2001.

\bibitem{ames2017}
A.~D. Ames, X.~Xu, J.~W. Grizzle, and P.~Tabuada, ``Control barrier function
based quadratic programs for safety critical systems,'' \emph{IEEE Trans.
Autom. Control}, vol.~62, no.~8, pp. 3861--3876, 2017.

\bibitem{smith1957}
O.~J.~M. Smith, ``Closer control of loops with dead time,'' \emph{Chem. Eng.
Prog.}, vol.~53, no.~5, pp. 217--219, 1957.

\bibitem{manitius1979}
A.~Z. Manitius and A.~W. Olbrot, ``Finite spectrum assignment problem for
systems with delays,'' \emph{IEEE Trans. Autom. Control}, vol.~24, no.~4,
pp. 541--552, 1979.

\bibitem{krstic2009}
M.~Krsti\'{c}, \emph{Delay Compensation for Nonlinear, Adaptive, and PDE
Systems}.\hskip 1em plus 0.5em minus 0.4em\relax Boston, MA: Birkh{\"a}user, 2009.

\bibitem{hale1993}
J.~K. Hale and S.~M. Verduyn Lunel, \emph{Introduction to Functional
Differential Equations}.\hskip 1em plus 0.5em minus 0.4em\relax New York, NY:
Springer-Verlag, 1993.

\bibitem{michiels2007}
W.~Michiels and S.-I. Niculescu, \emph{Stability and Stabilization of Time-Delay
Systems: An Eigenvalue-Based Approach}.\hskip 1em plus 0.5em minus 0.4em\relax
Philadelphia, PA: SIAM, 2007.

\bibitem{mayne2014}
D.~Q. Mayne, ``Model predictive control: Recent developments and future
promise,'' \emph{Automatica}, vol.~50, no.~12, pp. 2967--2986, 2014.

\bibitem{iea2024}
International Energy Agency, \emph{Electricity 2024: Analysis and Forecast to
2026}.\hskip 1em plus 0.5em minus 0.4em\relax Paris, France: IEA, 2024.

\end{thebibliography}

\end{document}